\uuid{qwSb}
\titre{Primitives usuelles}
\auteur{hanani ; mzari}
\organisation{exo7}
\licence{CC BY-NC-SA 4.0}
\datecreate{2020-01-20}
\niveau{L1}
\module{Analyse}
\chapitre{Calcul d'intégrales}
\sousChapitre{Primitives diverses}
\difficulte{2}
\theme{Primitives, primitives usuelles}
\exo7qcmid{lille-2.primitives.018}

\contenu{
\begin{question}
\qclassification{127.11}
\qtags{motcle=primitives usuelles}
Parmi les affirmations suivantes, cocher celles qui sont vraies :
\begin{answers}
\bad{$\displaystyle \int x\sqrt{x^2+1}\, \mathrm{d}x=\frac{x^2}{2}(1+x^2)^{3/2}+k$, $k\in \Rr$.}
\good{$\displaystyle \int x\sqrt{x^2+1}\, \mathrm{d}x=\frac{1}{3}(1+x^2)^{3/2}+k$, $k\in \Rr$.}
\bad{$\displaystyle \int \frac{x\, \mathrm{d}x}{\sqrt{x^2+1}}=\frac{x^2}{2}\frac{1}{\sqrt{x^2+1}}+k$, $k\in \Rr$.}
\good{$\displaystyle \int \frac{x\, \mathrm{d}x}{\sqrt{x^2+1}}=\sqrt{x^2+1}+k$, $k\in \Rr$.}
\end{answers}
\begin{explanations}
Avec $u=x^2+1$, on a $\mathrm{d}u=2x\mathrm{d}x$. D'où
$$\int x\sqrt{x^2+1}\mathrm{d}x=\frac{1}{2}\int \sqrt{u}\mathrm{d}u=\frac{u^{3/2}}{3}+k=\frac{1}{3}(1+x^2)^{3/2}+k,\; k\in \Rr,$$
et
$$\int \frac{x\mathrm{d}x}{\sqrt{x^2+1}}=\int \frac{\mathrm{d}u}{2\sqrt{u}}=\sqrt{u}+k=\sqrt{x^2+1}+k,\; k\in \Rr.$$
\end{explanations}
\end{question}
}
